\documentclass[11pt]{article}
\usepackage{amsmath,amssymb}
\title{Room 4 (P4 D-algebraicity): moderator summary}
\date{2026-09-26}
\begin{document}
\maketitle

Two seats (Tao, Erd\H os), both UNREVIEWED at the seat level but with extra care taken given that
this exact question already produced one retracted claim (\texttt{P4\_v2\_note.tex}).

\textbf{Tao's redo of the pole-rate summation: BOUND STILL ONLY LINEAR-TYPE, this time correctly
derived.} Using the actual general formula from \texttt{P1f\_lemma.tex} ($\Delta V=
(\pi is/N)(L_n+i\pi s/N)$, $u_j$ spaced $2\pi/|\Delta V|$ apart) together with the uniform bound
$|L_\mu|\le1.727$ for all $N\ge16$ (\texttt{lem:LS}(b)), the correct scaling is $|\Delta V|=
\Theta(1/N)$, hence $1-|q^\ast_j(\zeta)|=\Theta(1/(jN))$: the per-$\zeta$ coefficient \emph{shrinks}
like $1/N$, not the constant-in-$N$ or growing-in-$N$ behaviour either the original retracted
attempt or this room's working hypothesis considered. Summing over $N\le M$ gives an aggregate
coefficient of only $\Theta(M)$ -- still linear growth in the cutoff, not superpolynomial, and two
further gaps (non-uniform $j_0(\zeta)$, non-uniform ``$r$ large enough'' threshold) would likely
weaken this further. This \emph{confirms}, with an actual derivation, what the earlier retraction's
correction paragraph asserted but did not derive -- no new false claim, no progress either.

\textbf{Erd\H os's alternative-obstruction search: NOTHING WORKED, precisely diagnosed.} Three
elementary candidates (functional-equation degree incompatibility, Wronskian/algebraic-independence
search, Cauchy--Kovalevskaya parameter counting) were checked. Each either needs exactly the missing
quantitative growth/counting ingredient \texttt{P4\_note.tex} already names, or is structurally
blocked by results paper2a already proves in the opposite direction (\texttt{prop:noqdiff},
\texttt{cor:Qminusone} rule out the functional relation a degree-counting argument would need).

\textbf{Room verdict: NO PROGRESS,} honestly. Both seats independently converge on the same
conclusion: the missing ingredient (a quantitative, superpolynomial-strength pole-counting or growth
bound) is not reachable by the routes tried, and no alternative obstruction was found either. This
is not a failure of the room -- it correctly avoided repeating the earlier error, and it closes off
two more candidate routes with precise reasons rather than leaving them dangling. P4 stands exactly
where \texttt{P4\_note.tex} (v10.19.0) left it: OPEN, with the right theorem family named and the
same two missing ingredients.

\end{document}
